\documentclass[sigconf,anonymous]{acmart}

%% Remove ACM copyright/DOI for submission
\setcopyright{none}
\settopmatter{printacmref=false}
\renewcommand\footnotetextcopyrightpermission[1]{}
\pagestyle{plain}

\usepackage{booktabs}
\usepackage{subcaption}
\usepackage{xspace}
\usepackage{enumitem}

%% Convenience macros — update these as data collection continues
\newcommand{\sys}{EduThreat-CTI\xspace}
\newcommand{\rawcollected}{1{,}360\xspace}       % Total ingested incidents (after noise removal)
\newcommand{\numincidents}{845\xspace}            % Verified education cyber incidents (is_education_related=1)
\newcommand{\numenriched}{891\xspace}             % Incidents processed through LLM enrichment
\newcommand{\numpending}{469\xspace}              % Awaiting enrichment
\newcommand{\numcountries}{47\xspace}
\newcommand{\numsources}{15\xspace}
\newcommand{\numransomware}{200\xspace}
\newcommand{\numdatabreach}{425\xspace}
\newcommand{\numactors}{134\xspace}
\newcommand{\numfamilies}{49\xspace}
\newcommand{\avgrecovery}{27.6\xspace}
\newcommand{\totaldemanded}{\$346M\xspace}
\newcommand{\paymentrate}{22.2\%\xspace}

\begin{document}

\title{EduThreat-CTI: An Automated Cyber Threat Intelligence Pipeline\\for the Education Sector}

\author{Anonymous Author(s)}
\affiliation{%
  \institution{Anonymous Institution}
}

\begin{abstract}
Academic institutions are increasingly targeted by cyber adversaries, yet the education sector lacks dedicated, large-scale cyber threat intelligence (CTI) infrastructure comparable to that available for healthcare or finance. We present \sys, the first open-source, automated CTI pipeline purpose-built for the education sector. \sys continuously ingests reports from \numsources diverse OSINT sources---including curated databases, ransomware trackers, security news outlets, and RSS feeds---processes \rawcollected candidate incidents through cross-source deduplication and LLM-based relevance verification, and applies schema-constrained enrichment to produce structured records with 192+ CTI fields per incident, including MITRE ATT\&CK mappings, attack timelines, and multidimensional impact metrics.

We deploy \sys over a 9-year collection window (2017--2026) and present what is, to our knowledge, the largest measurement study of cyber incidents targeting educational institutions. Our pipeline identifies and verifies \numincidents education-sector cyber incidents across \numcountries countries. Key findings include: (1)~ransomware (encryption and double extortion combined) accounts for 24.0\% of categorized incidents, with Vice Society, Clop, and Medusa among the most active families; (2)~the United States absorbs 68.9\% of reported incidents, revealing extreme geographic concentration; (3)~a transparency crisis persists, with 29.5\% of incidents lacking any identified initial access vector; and (4)~ransom economics across \numransomware ransomware incidents show \totaldemanded in total demands against a \paymentrate payment rate. We release \sys, the enriched dataset, and a production analytics dashboard as open resources for the research community.
\end{abstract}

\begin{CCSCLASS}
\ccsdesc[500]{Security and privacy~Intrusion/anomaly detection and malware mitigation}
\ccsdesc[300]{Security and privacy~Security services}
\ccsdesc[300]{Information systems~Data mining}
\end{CCSCLASS}

\keywords{cyber threat intelligence, education sector, ransomware, OSINT, LLM enrichment, measurement study, MITRE ATT\&CK}

\maketitle

%% ============================================================
%% SECTION 1: INTRODUCTION
%% ============================================================
\section{Introduction}\label{sec:intro}

Academic institutions occupy a uniquely precarious position in the global cybersecurity landscape. Universities and research institutes house vast repositories of sensitive data---student records, intellectual property, federally funded research, and healthcare information---while operating inherently open networks designed to foster collaboration~\cite{garrison2011longitudinal}. This openness, combined with decentralized IT governance, constrained budgets, and a heterogeneous user population spanning students, faculty, and administrative staff, creates attack surfaces that adversaries readily exploit.

The scale of the problem is substantial. Industry reports document a 70\% year-over-year increase in ransomware attacks against higher education in 2023 alone~\cite{malwarebytes2024education}, with the Sophos State of Ransomware in Education survey finding that 66\% of higher-education organizations were hit by ransomware in the preceding year and that the median ransom payment reached \$4.4 million~\cite{sophos2024education}. Yet despite this escalating threat, the education sector remains conspicuously underserved by the cyber threat intelligence (CTI) community. Existing CTI platforms---MISP~\cite{wagner2016misp}, OpenCTI~\cite{opencti}, and commercial feeds---are sector-agnostic, requiring analysts to manually filter for education-relevant incidents. No dedicated, continuously updated, open-source CTI resource exists for academia.

This gap has three consequences. \textit{First}, institutions lack sector-specific situational awareness: without a consolidated view of which threat actors target universities, through which vectors, and with what operational impact, individual institutions cannot benchmark their defenses. \textit{Second}, researchers studying the education threat landscape are forced to assemble ad-hoc datasets from heterogeneous sources, limiting reproducibility and cross-study comparability. \textit{Third}, the transparency deficit---where institutions suppress disclosure for reputational reasons~\cite{sophos2024education}---is compounded by the absence of infrastructure that could aggregate and anonymize incident data at scale.

We address these gaps with \sys, the first open-source, automated CTI pipeline purpose-built for the education sector. \sys implements a three-phase architecture: (1)~\textit{Ingestion}, which continuously collects candidate incidents from \numsources diverse OSINT sources---curated databases, ransomware trackers, security news, and RSS feeds---using source-specific scrapers, incremental state tracking, and cross-source URL-based deduplication; (2)~\textit{Enrichment}, which retrieves full-text articles through a multi-fallback pipeline (newspaper3k $\rightarrow$ Selenium $\rightarrow$ archive.org), performs LLM-based education-relevance verification to filter noise from keyword-matched scraping, and applies schema-constrained extraction to produce 192+ structured CTI fields per verified incident; and (3)~\textit{Analytics}, which serves the enriched corpus through a REST API with 40+ analytical endpoints and an interactive dashboard.

Using \sys, we conduct what is, to our knowledge, the largest systematic OSINT-based measurement study of cyber incidents targeting educational institutions. Over a 9-year window (2017--2026), our pipeline ingests \rawcollected candidate incidents from \numsources sources, from which LLM-assisted verification confirms \numincidents education-sector cyber incidents across \numcountries countries---each enriched with comprehensive CTI metadata spanning 192+ structured fields. Unlike prior work that relies on single curated sources, annual surveys, or qualitative reviews, \sys automates the full intelligence cycle from multi-source collection through structured enrichment to analytical serving, and runs continuously as new incidents are reported.

\smallskip
\noindent\textbf{Contributions.} This paper makes four contributions:

\begin{enumerate}[leftmargin=*,nosep]
\item \textbf{System.} We design and implement \sys, a production-grade, open-source CTI pipeline that automates the full cycle from OSINT collection through LLM-assisted enrichment to analytical serving, tailored specifically for the education sector (\S\ref{sec:system}).

\item \textbf{Dataset.} We construct and release the largest education-sector cyber incident corpus to date: \numincidents verified incidents spanning \numcountries countries and 9 years, each enriched to 192+ structured CTI fields including MITRE ATT\&CK mappings, attack timelines, and impact metrics---identified from \rawcollected candidates through automated relevance verification (\S\ref{sec:dataset}).

\item \textbf{Measurement.} We present a comprehensive empirical characterization of the education threat landscape, revealing ransomware dominance, extreme geographic concentration, systemic supply-chain exposure, and a pervasive transparency crisis (\S\ref{sec:measurement}).

\item \textbf{Open resource.} We release the pipeline code, enriched dataset, REST API, and interactive dashboard as open resources to support future research and institutional decision-making.
\end{enumerate}


%% ============================================================
%% SECTION 2: BACKGROUND AND RELATED WORK
%% ============================================================
\section{Background \& Related Work}\label{sec:related}

We situate our work at the intersection of three research threads: (1)~cybersecurity in the education sector, (2)~cyber threat intelligence platforms and pipelines, and (3)~LLM-assisted information extraction for security.

%% --- 2.1 ---
\subsection{Cybersecurity in the Education Sector}\label{sec:related:edu}

The education sector has been recognized as a high-risk target for over a decade. Garrison and Ncube~\cite{garrison2011longitudinal} documented 290 university breaches exposing 10.8 million records between 2005 and 2009, identifying stolen devices and hacking as primary causes. More recently, industry measurement efforts have quantified the acceleration: Malwarebytes reported a 70\% year-over-year increase in higher-education ransomware (68 attacks in 2022 vs.\ 116 in 2023)~\cite{malwarebytes2024education}, while Sophos found that 66\% of surveyed higher-education organizations experienced a ransomware incident in the prior year, with median ransom payments of \$4.4M and recovery costs averaging \$4.02M~\cite{sophos2024education}. Comparitech estimated that ransomware attacks on U.S.\ schools and colleges have cost \$9.45 billion cumulatively~\cite{comparitech2024schools}.

However, these reports suffer from three shared limitations. First, they are either narrow in geographic scope (e.g., single-country reviews~\cite{lallie2025ukuniversities}), restricted to a single attack type, or rely on annual survey self-reports rather than systematic OSINT collection~\cite{sophos2024education,malwarebytes2024education}. Second, existing academic studies are largely qualitative: Lallie et al.~\cite{lallie2025ukuniversities} provide a chronological narrative of UK university attacks, and Jawaid~\cite{jawaid2022eduthreats} surveys institutional threats without an empirical corpus. Third, no prior work provides an open, reproducible, continuously updated dataset enriched with structured CTI metadata. \sys addresses all three gaps by combining automated multi-source ingestion from \numsources OSINT sources with LLM-driven relevance verification and 192-field enrichment across \numcountries countries, yielding \numincidents verified incidents from \rawcollected candidates.

%% --- 2.2 ---
\subsection{Cyber Threat Intelligence Platforms}\label{sec:related:cti}

The CTI ecosystem spans both open-source and commercial platforms. MISP (Malware Information Sharing Platform)~\cite{wagner2016misp} enables collaborative sharing and correlation of indicators of compromise (IoCs) through a structured event model. OpenCTI~\cite{opencti} provides a graph-based knowledge management platform aligned with STIX/TAXII standards. On the commercial side, platforms such as Recorded Future, Mandiant, and CrowdStrike aggregate threat feeds from diverse sources.

A parallel research thread has explored automated CTI pipelines. Sun et al.~\cite{sun2023ctisurvey} provide a comprehensive survey of CTI mining techniques for proactive defense, while Furumoto et al.~\cite{furumoto2025ctisurvey} present a measurement-based analysis of over 200 CTI studies spanning 2001--2025. On the pipeline side, Tundis et al.~\cite{tundis2022feature} developed a feature-driven method for automated assessment of OSINT CTI source quality, proposing scoring functions to quantify relevance.

These platforms and pipelines are \textit{sector-agnostic}: they neither target the education sector specifically nor provide the domain-specific schema, enrichment, or analytics that education-sector CTI requires. \sys fills this gap by combining multi-source OSINT ingestion with an education-tailored enrichment schema spanning 192+ fields, including institution type, operational impact on teaching and research, and education-specific regulatory exposure (FERPA, GDPR).

%% --- 2.3 ---
\subsection{LLM-Assisted CTI Extraction}\label{sec:related:llm}

Recent work has explored large language models for automated CTI extraction. Rani et al.~\cite{rani2024ttpxhunter} developed TTPXHunter, which uses SecureBERT to extract TTPs from threat reports in STIX format. Wang et al.~\cite{wang2024knowcti} proposed KnowCTI, incorporating cybersecurity domain knowledge to enhance entity and relation extraction from CTI text. Fieblinger et al.~\cite{fieblinger2024kgllmcti} combined knowledge graphs with LLMs for actionable CTI extraction, evaluating prompt engineering and fine-tuning approaches. More broadly, Jaffal et al.~\cite{jaffal2025llmcybersec} surveyed 223 studies on LLM applications in cybersecurity published between 2021 and 2025.

A critical limitation identified by Mezzi et al.~\cite{mezzi2025unreliable} is that LLMs show strong performance on short inputs but degrade on realistic full-length reports (averaging 3{,}009 words), raising concerns about reliability for CTI extraction at scale. Nguyen et al.~\cite{nguyen2025mitreextract} further demonstrated challenges in automated MITRE ATT\&CK technique identification from CTI reports.

Our approach addresses these challenges through schema-constrained extraction: rather than open-ended generation, we provide the LLM with a detailed JSON schema (192+ fields with controlled enumerations) and validate outputs against Pydantic models, achieving structured extraction at scale while mitigating hallucination through type enforcement and enumeration constraints. We evaluate extraction quality through manual validation of sampled records (\S\ref{sec:dataset}).


%% ============================================================
%% SECTION 3: CONTRIBUTIONS & NOVELTY
%% ============================================================
\section{Contributions \& Novelty}\label{sec:contributions}

\sys makes four distinct contributions that, taken together, address the absence of dedicated CTI infrastructure for the education sector. Table~\ref{tab:comparison} positions each contribution against representative prior work.

%% --- 3.1 ---
\subsection{An Automated, Education-Specific CTI Pipeline}\label{sec:contrib:system}

Existing open-source CTI platforms---MISP~\cite{wagner2016misp} and OpenCTI~\cite{opencti}---are sector-agnostic ingestion and sharing frameworks. They require manual feed curation, do not perform education-relevance filtering, and provide no domain-tailored enrichment schema. Commercial feeds (Recorded Future, Mandiant) are similarly generic and closed.

\sys is the first CTI pipeline designed ground-up for the education sector. Its novel engineering contributions are:

\begin{itemize}[leftmargin=*,nosep]
  \item \textbf{Heterogeneous multi-source ingestion.} Fifteen source-specific scrapers cover curated databases (KonBriefing~\cite{konbriefing}), ransomware leak-site mirrors (Ransomware.live~\cite{ransomwarelive}), breach archives (DataBreaches.net~\cite{databreach2024}), five security news outlets, and five RSS feeds. Each scraper implements incremental state tracking so that only new incidents are processed on each run.

  \item \textbf{Cross-source deduplication.} A URL-normalization and hash-based merging layer deduplicates across all \numsources sources, preserving multi-source attribution while emitting each incident exactly once. This is non-trivial: the same incident is routinely reported by 3--6 outlets with divergent URLs, titles, and dates.

  \item \textbf{Multi-fallback article retrieval.} Full-text fetching uses a three-tier fallback: direct HTTP retrieval via newspaper3k, JavaScript-rendered fetching via Playwright/Selenium for bot-protected sources, and archive.org as a last resort for expired pages. This maximises corpus completeness across source types.

  \item \textbf{Schema-constrained LLM enrichment.} A structured JSON schema (192+ typed fields with controlled enumerations) is provided to a large language model at inference time. All outputs are validated against Pydantic models; type violations and out-of-vocabulary values are rejected and flagged for re-processing. This enforces structure at scale and directly addresses the hallucination degradation documented for open-ended CTI extraction~\cite{mezzi2025unreliable}.

  \item \textbf{Continuous production operation.} \sys runs as a scheduled service on Railway, re-ingesting new incidents without human intervention. The pipeline has been in continuous operation since deployment, with the corpus growing automatically as new incidents are reported.
\end{itemize}

%% --- 3.2 ---
\subsection{The EduThreat Corpus: A Structured Education CTI Dataset}\label{sec:contrib:dataset}

No prior open dataset provides structured, multi-field CTI records for education-sector cyber incidents at scale. Existing resources are either raw unstructured lists (KonBriefing), paywalled surveys (Sophos, Malwarebytes), or single-country qualitative reviews (Lallie et al.~\cite{lallie2025ukuniversities}).

The EduThreat corpus comprises \numincidents verified education-sector incidents across \numcountries countries, spanning 2017--2026, and is enriched with 192+ structured CTI fields per incident. Fields novel to this corpus include:

\begin{itemize}[leftmargin=*,nosep]
  \item \textbf{Education-specific institution taxonomy:} university, college, K-12 school, research institute, hospital-affiliated, consortium.
  \item \textbf{Operational impact on core functions:} teaching disruption, research slowdown, admissions suspension, library/HPC outage---categories absent from generic CTI schemas.
  \item \textbf{Regulatory exposure flags:} FERPA, GDPR, HIPAA, and state-level breach-notification obligations per incident.
  \item \textbf{Transparency scoring:} disclosure timing, detection source, update frequency, and post-incident communication quality---enabling empirical study of institutional disclosure behaviour.
  \item \textbf{MITRE ATT\&CK mapping:} tactic and technique labels extracted per incident, enabling cross-incident TTP analysis without requiring manual annotation.
  \item \textbf{Ransom economics:} demanded amounts, paid amounts, cryptocurrency denomination, and negotiation outcome where disclosed.
\end{itemize}

Table~\ref{tab:comparison} compares the EduThreat corpus against representative prior resources along key dimensions.

\begin{table}[t]
\centering
\caption{Comparison of education-sector cybersecurity resources. ``Automated'' = pipeline runs without manual intervention per incident. ``Open'' = freely available without registration or payment.}
\label{tab:comparison}
\small
\begin{tabular}{lcccc}
\toprule
& \textbf{KonBriefing} & \textbf{Sophos /} & \textbf{Lallie} & \textbf{\sys} \\
& \textbf{(raw list)} & \textbf{Malwarebytes} & \textbf{et al.~\cite{lallie2025ukuniversities}} & \textbf{(ours)} \\
\midrule
Incidents        & $\sim$600   & Survey-based  & Qualitative   & \numincidents \\
Countries        & 40+         & 14            & UK only       & \numcountries \\
CTI fields       & $\sim$5     & $\sim$8       & None          & \textbf{192+} \\
MITRE ATT\&CK   & \texttimes  & \texttimes    & \texttimes    & \checkmark    \\
Automated        & \texttimes  & Annual survey & \texttimes    & \checkmark    \\
Open             & \checkmark  & \texttimes    & \checkmark    & \checkmark    \\
API + Dashboard  & \texttimes  & \texttimes    & \texttimes    & \checkmark    \\
Continuously updated & \texttimes & \texttimes & \texttimes   & \checkmark    \\
\bottomrule
\end{tabular}
\end{table}

%% --- 3.3 ---
\subsection{A Large-Scale Measurement of the Education Threat Landscape}\label{sec:contrib:measurement}

Leveraging the EduThreat corpus, we present the most comprehensive empirical characterisation of education-sector cyber threats to date. Prior measurement work in adjacent sectors (healthcare~\cite{sophos2024education}, government) relies on surveys or single-source datasets. Our measurement study reveals several findings not previously documented at this scale:

\begin{itemize}[leftmargin=*,nosep]
  \item \textbf{Geographic inequality:} Incident distribution across \numcountries countries follows a power-law, with the United States accounting for 68.9\% of reported incidents. We quantify this concentration empirically rather than via survey approximation.

  \item \textbf{Ransomware dominance and economics:} Ransomware (encryption + double extortion) accounts for 24.0\% of categorised incidents. Across \numransomware ransomware incidents, we document \totaldemanded in total demands against a \paymentrate payment rate---the first such estimate derived from a systematic OSINT corpus for education.

  \item \textbf{Transparency crisis:} 29.5\% of incidents have no identified initial access vector; disclosure-timeline data is missing for the majority of cases. We characterise this structural opacity and its implications for sector-wide learning.

  \item \textbf{Threat actor and TTP profiling:} \numactors distinct threat actors and \numfamilies ransomware families are identified across the corpus, with MITRE ATT\&CK techniques mapped for a representative subset, enabling the first TTP-level view of education-targeting adversaries.

  \item \textbf{Operational impact:} Average recovery time is \avgrecovery days, with ransomware-induced outages following a heavy-tailed distribution (median 72h, P90 $\approx$ 20 days).
\end{itemize}

%% --- 3.4 ---
\subsection{Open Science Artifact}\label{sec:contrib:openscience}

In alignment with ACM's open science policy, we release all artifacts associated with this paper:

\begin{itemize}[leftmargin=*,nosep]
  \item \textbf{Pipeline code:} Full source code for all three phases (ingestion, enrichment, analytics) released under an open-source licence. [Anonymous repository provided at submission.]

  \item \textbf{Enriched dataset:} The EduThreat corpus in CSV and JSON formats, with identifiable fields anonymised where required by institutional privacy obligations.

  \item \textbf{REST API:} A production REST API exposing 40+ analytical endpoints over the live corpus, documented at [anonymous URL].

  \item \textbf{Interactive dashboard:} A web-based analytics dashboard providing geographic, temporal, actor, and TTP views of the corpus without requiring API access.
\end{itemize}

%% ============================================================
%% Placeholder sections (to be written)
%% ============================================================
\section{System Design}\label{sec:system}
% To be written next

\section{Dataset \& Methodology}\label{sec:dataset}
% To be written

\section{Measurement Study}\label{sec:measurement}
% To be written

\section{Discussion}\label{sec:discussion}
% To be written

\section{Conclusion}\label{sec:conclusion}
% To be written


%% ============================================================
%% REFERENCES
%% ============================================================
\bibliographystyle{ACM-Reference-Format}
\bibliography{references}

\end{document}
